% Pollis Domain Packages: Heteroskedasticity and AR(1) errors (HARE.jl) -- draft 1 (28 October 2026)
% Compile with: latexmk -pdf HARE.tex && latexmk -c
% Shared look for the Pollis New Methods documents (same fonts and colours as the brochure and the workflow handouts).
\documentclass[10pt,a4paper]{article}

\usepackage[a4paper,left=18mm,right=18mm,top=15mm,bottom=17mm,footskip=8mm]{geometry}
\usepackage[T1]{fontenc}
\usepackage[default,semibold]{sourcesanspro}
\usepackage{textcomp}
\usepackage[italic]{mathastext}
\usepackage{amsmath}
\usepackage{microtype}
\usepackage[table]{xcolor}
\usepackage{tikz}
\usepackage[most]{tcolorbox}
\usepackage{fontawesome5}
\usepackage{tabularx,booktabs}
\usepackage{enumitem}
\usepackage{fancyhdr}
\usepackage[hidelinks]{hyperref}

% ---------------------------------------------------------------- colours (as in the brochure)
\definecolor{navy}{HTML}{1C3A5E}
\definecolor{pblue}{HTML}{3B78B5}
\definecolor{amber}{HTML}{F2B33D}
\definecolor{deepamber}{HTML}{D9902A}
\definecolor{pale}{HTML}{EEF3F9}
\definecolor{ink}{HTML}{2A2F36}
\definecolor{logonavy}{HTML}{173D63}
\definecolor{logodark}{HTML}{C98A1E}
\definecolor{logoamber}{HTML}{F2B446}
\definecolor{logocream}{HTML}{FEEFD2}
\definecolor{soft}{HTML}{5E6773}
\definecolor{agree}{HTML}{2F7D4F}
\definecolor{differ}{HTML}{B23B3B}
\color{ink}

\setlength{\parindent}{0pt}
\setlength{\parskip}{4pt}
\setlist{nosep,leftmargin=*,itemsep=2pt}
\renewcommand{\arraystretch}{1.15}

\pagestyle{fancy}
\fancyhf{}
\renewcommand{\headrulewidth}{0pt}
\fancyfoot[L]{\footnotesize\color{soft}\textbf{Pollis}\enspace New Methods\enspace·\enspace\nmtheme}
\fancyfoot[R]{\footnotesize\color{soft}\thepage}
\newcommand{\nmtheme}{}
\newcommand{\nmdate}{11 October 2026}

% The Pollis logo, redrawn so it prints sharply at any size
\newcommand{\pollislogo}[1]{%
  \begin{tikzpicture}[x=#1/512,y=-#1/512,baseline=(current bounding box.center)]
    \fill[logonavy] (0,0) rectangle (512,512);
    \fill[logodark] (36,422) -- (104,256) -- (172,422) -- cycle;
    \fill[logoamber] (104,422) -- (203,180) -- (302,422) -- cycle;
    \fill[logocream] (203,422) -- (339,89) -- (476,422) -- cycle;
  \end{tikzpicture}}

% Top of the first page: logo, name; then the kicker, title and a one-line promise
\newcommand{\nmhead}[3]{%
  \renewcommand{\nmtheme}{#1}%
  \thispagestyle{fancy}%
  \begin{minipage}[c]{13mm}\pollislogo{11mm}\end{minipage}%
  \begin{minipage}[c]{\dimexpr\linewidth-13mm\relax}
    {\color{navy}\Large\bfseries Pollis}\enspace{\color{soft}New Methods}\hfill
    {\color{deepamber}\small\bfseries\MakeUppercase{#1}}
  \end{minipage}\par\vspace{7mm}
  {\color{navy}\fontsize{22}{26}\selectfont\bfseries #2}\par\vspace{4pt}
  {\color{soft}\large #3}\par\vspace{5mm}}

\newcommand{\nmsection}[1]{\par\vspace{5pt}{\color{navy}\large\bfseries #1}\par\vspace{1pt}}
\newcommand{\nmsub}[1]{\par\vspace{2pt}{\color{deepamber}\bfseries #1}\par}
\newcommand{\gt}{\,{\color{deepamber}\bfseries ›}\,}
\newcommand{\menupath}[1]{{\bfseries\color{navy}#1}}
\newcommand{\fn}[1]{\texttt{\small #1}}
% Citations: \nmcite{key}{Breiman (2001)} links to the entry \nmref{key}{...} in the References
\newcommand{\nmcite}[2]{\hyperlink{ref:#1}{\color{pblue}#2}}
\newcommand{\nmref}[2]{\raisebox{1.1em}[0pt][0pt]{\hypertarget{ref:#1}{}}#2\par}
\newcommand{\nmdoi}[1]{\href{https://doi.org/#1}{\color{pblue}doi:#1}}
% Where a notebook is in Pollis: its title in Browse Methods, and its file
\newcommand{\browse}[2]{\menupath{Compose \gt New Methods \gt Browse Methods}: \emph{#1} (\fn{#2.ipynb})}

% Key facts at the top: source, missing in Julia, notebook, result
\newcommand{\factsheet}[4]{%
  \begin{tcolorbox}[enhanced, colback=pale, colframe=pale, boxrule=0pt, sharp corners,
    left=7pt, right=7pt, top=5pt, bottom=5pt, fontupper=\small]
  \begin{tabularx}{\linewidth}{@{}>{\leavevmode\bfseries\color{navy}}p{27mm}>{\leavevmode}X@{}}
  Tested against & #1\\
  Missing in Julia & #2\\
  Notebook & #3\\
  Result & #4
  \end{tabularx}
  \end{tcolorbox}}

% Boxes
\newtcolorbox{pbox}[2][]{enhanced, colback=pale, colframe=amber, boxrule=0pt,
  leftrule=3pt, sharp corners, left=7pt, right=7pt, top=6pt, bottom=6pt,
  fontupper=\small\raggedright, before upper={{\color{navy}\normalsize\bfseries #2}\par\smallskip}, #1}
\newtcolorbox{nbox}[2][]{enhanced, colback=navy, colframe=navy, boxrule=0pt,
  sharp corners, left=8pt, right=8pt, top=7pt, bottom=7pt, coltext=white,
  fontupper=\small\raggedright, before upper={{\color{amber}\normalsize\bfseries #2}\par\smallskip}, #1}

% Value tables: published beside Julia
\newenvironment{valuetable}[1]{%
  \par\small\tabularx{\linewidth}{#1}\toprule}{\bottomrule\endtabularx\par}
\newcommand{\ok}{{\color{agree}\faCheck}}
\newcommand{\no}{{\color{differ}\faTimes}}

% End of the document: credit and licence
\newcommand{\nmcredit}{%
  \par\vfill
  \begin{minipage}[b]{0.78\linewidth}\footnotesize\color{soft}
  Methods written and tested with Claude Code (Anthropic) in Pollis, from the published papers and their data.\par
  \textcopyright\ 2026 Antonio Saragga Seabra\enspace·\enspace Licensed under \href{https://creativecommons.org/licenses/by-sa/4.0/}{CC BY-SA 4.0}\enspace·\enspace Draft, \nmdate
  \end{minipage}\hfill
  \begin{minipage}[b]{0.18\linewidth}\raggedleft\pollislogo{9mm}\par\vspace{2pt}{\footnotesize contact@pollis.org}\end{minipage}}


% A package Pollis needs for panels of its own, not one that reproduces a paper: its own header, key facts and credit
\renewcommand{\nmdate}{28 October 2026}
\fancyfoot[L]{\footnotesize\color{soft}\textbf{Pollis}\enspace Domain Packages\enspace·\enspace\nmtheme}
\renewcommand{\nmhead}[3]{%
  \renewcommand{\nmtheme}{#1}%
  \thispagestyle{fancy}%
  \begin{minipage}[c]{13mm}\pollislogo{11mm}\end{minipage}%
  \begin{minipage}[c]{\dimexpr\linewidth-13mm\relax}
    {\color{navy}\Large\bfseries Pollis}\enspace{\color{soft}Domain Packages}\hfill
    {\color{deepamber}\small\bfseries\MakeUppercase{#1}}
  \end{minipage}\par\vspace{7mm}
  {\color{navy}\fontsize{22}{26}\selectfont\bfseries #2}\par\vspace{4pt}
  {\color{soft}\large #3}\par\vspace{5mm}}
\renewcommand{\factsheet}[4]{%
  \begin{tcolorbox}[enhanced, colback=pale, colframe=pale, boxrule=0pt, sharp corners,
    left=7pt, right=7pt, top=5pt, bottom=5pt, fontupper=\small]
  \begin{tabularx}{\linewidth}{@{}>{\leavevmode\bfseries\color{navy}}p{27mm}>{\leavevmode}X@{}}
  In Pollis & #1\\
  Missing in Julia & #2\\
  Package & #3\\
  Tests & #4
  \end{tabularx}
  \end{tcolorbox}}
\renewcommand{\nmcredit}{%
  \par\vfill
  \begin{minipage}[b]{0.78\linewidth}\footnotesize\color{soft}
  Package written and tested with Claude Code (Anthropic) in Pollis.\par
  \textcopyright\ 2026 Antonio Saragga Seabra\enspace·\enspace Licensed under \href{https://creativecommons.org/licenses/by-sa/4.0/}{CC BY-SA 4.0}\enspace·\enspace Draft, \nmdate
  \end{minipage}\hfill
  \begin{minipage}[b]{0.18\linewidth}\raggedleft\pollislogo{9mm}\par\vspace{2pt}{\footnotesize contact@pollis.org}\end{minipage}}

\begin{document}
\emergencystretch=3em

\nmhead{Linear models}
  {HARE: errors that change size, errors that remember}
  {A regression whose errors grow with a regressor, or carry over from one period to the next, or both: how to
   estimate it well, and how to tell which it is.}

\factsheet
  {\menupath{Model \gt Linear Models}: \emph{Linear Models with Heteroskedasticity}, \emph{Linear Models with
   Autocorrelation} and \emph{Linear Models with Heteroskedasticity + AR(1)}. Each panel's Code Example runs the
   package's estimators and tests on its data.}
  {feasible GLS with a model of the variance (exponential, power, quadratic, linear, by group), by two-step or
   iterated weighted least squares and by maximum likelihood; the AR(1) estimators of Prais and Winsten, Cochrane and
   Orcutt, Hildreth and Lu, and Beach and MacKinnon; the two together; Harvey's and Glejser's tests. GLM.jl weighs
   by weights you give it, and robust standard errors leave the estimates as they are.}
  {\fn{HARE.jl} (Heteroskedasticity and Autocorrelation Estimators), from GitHub:
   \fn{Pkg.add(url = "https://github.com/Trumpingtons/HARE.jl")}.}
  {1,295 tests of every estimator and test; 30 against R (prais, orcutt, skedastic, nlme) on the same data, to
   $10^{-8}$ or closer; 33 that recover the parameters from a million observations.}

\nmsection{Why a package}
Ordinary least squares stays unbiased when the errors are heteroskedastic or autocorrelated, but it is no longer the
most precise estimator, and its standard errors are wrong. Robust (White, Newey--West) standard errors fix the second
and leave the first. When the form of the variance or of the dependence is known, or can be estimated, generalized
least squares fixes both. In R, these estimators are spread over prais, orcutt, nlme and skedastic; in Stata over
\fn{prais} and \fn{hetregress}; in SAS most are in PROC AUTOREG. Julia had none of them, and Pollis' linear
model panels needed all of them in one place, with the same interface, results that print as GLM.jl's do, and the
estimate of the variance or of $\rho$ beside the coefficients.

\nmsub{Where they are used}
\begin{itemize}
  \item \textbf{Cross sections whose spread grows with size.} Household spending varies more among the rich, firms'
    costs more among large firms, regional outcomes more where the population is small: the variance is a function of
    a regressor, the setting of \nmcite{harvey1976}{Harvey (1976)} and \nmcite{glejser1969}{Glejser (1969)}.
  \item \textbf{Groups with variances of their own.} Countries, hospitals or survey waves measured with different
    precision: one variance a group (\nmcite{greene2018}{Greene 2018}, §9.7.2).
  \item \textbf{Time series with persistent errors.} Demand or a macroeconomic relation fitted to annual or quarterly
    data, where a shock lasts beyond its period: AR(1) errors (\nmcite{co1949}{Cochrane and Orcutt 1949};
    \nmcite{pw1954}{Prais and Winsten 1954}).
  \item \textbf{Both at once.} A series whose shocks persist and whose volatility moves with a regressor; the
    estimators of the third panel, for which we found no other implementation.
\end{itemize}

\nmsection{The method}
The model is $y_i = \mathbf{x}_i^\top\boldsymbol\beta + u_i$, and the package estimates $\boldsymbol\beta$ with a
model of $u_i$. Each function takes a formula and a table, \fn{f(@formula(y \textasciitilde{} x1 + x2), data; Z =
@formula(0 \textasciitilde{} x1))}, or matrices without a column of ones, \fn{f(X, y; Z)}; $\mathbf{z}_i$, the
regressors of the variance, are by default those of the mean.

\nmsub{Heteroskedasticity: a model of the variance}
\begin{valuetable}{@{}>{\raggedright\arraybackslash}p{24mm}>{\raggedright\arraybackslash}p{50mm}>{\raggedright\arraybackslash}X@{}}
\textbf{Link} & \textbf{Variance} & \textbf{Functions}\\\midrule
Exponential & $\sigma_i^2 = \exp(\gamma_0 + \mathbf{z}_i^\top\boldsymbol\gamma)$ \hfill Harvey & \fn{exponential\_fwls}, \fn{\_ifwls}, \fn{\_mle}\\
Power & $\log\sigma_i^2 = \gamma_0 + \sum_j \gamma_j \log|z_{ij}|$ & \fn{power\_fwls}, \fn{\_ifwls}, \fn{\_mle}\\
Quadratic & $\sigma_i = \gamma_0 + \mathbf{z}_i^\top\boldsymbol\gamma$ \hfill Glejser & \fn{quadratic\_fwls}, \fn{\_ifwls}, \fn{\_mle}\\
Linear & $\sigma_i^2 = \gamma_0 + \mathbf{z}_i^\top\boldsymbol\gamma$ \hfill Breusch--Pagan & \fn{linear\_fwls}, \fn{\_ifwls}, \fn{\_mle}\\
By group & $\sigma_i^2 = \sigma_g^2$ for $i$ in group $g$ & \fn{two\_step\_groupwise}, \fn{iterated\_groupwise}, \fn{groupwise\_mle}\\
\end{valuetable}
\emph{Two-step FWLS} fits OLS, regresses a function of the residuals on $\mathbf{z}_i$ (for the exponential link,
$\log\hat u_i^2$, less $\mathrm{E}\log\chi^2_1 = -1.2704$, so that every $\gamma$ is consistent and not only the
slopes), and refits by weighted least squares with weights $1/\hat\sigma_i^2$. It needs no optimizer and stays
consistent when the errors are not normal. \emph{Iterated FWLS} repeats the last two steps until $\boldsymbol\beta$
stops moving, and for the exponential link converges to the maximum likelihood estimate
(\nmcite{ok1974}{Oberhofer and Kmenta 1974}). \emph{MLE} maximizes the normal likelihood over $\boldsymbol\beta$
and $\boldsymbol\gamma$ together, by L-BFGS with ForwardDiff's gradients, started from the two-step estimate, and
gives standard errors for $\boldsymbol\gamma$ from the Hessian.

\nmsub{Autocorrelation: AR(1) errors, $u_t = \rho u_{t-1} + e_t$}
\begin{itemize}
  \item \textbf{Prais--Winsten} (\fn{two\_step\_prais\_winsten}, \fn{iterated\_prais\_winsten}): $\rho$ from the
    residuals' first-order regression, then GLS on $y_t - \rho y_{t-1}$, keeping the first observation scaled by
    $\sqrt{1-\rho^2}$.
  \item \textbf{Cochrane--Orcutt} (\fn{two\_step\_cochrane\_orcutt}, \fn{iterated\_cochrane\_orcutt}): the same,
    dropping the first observation.
  \item \textbf{Hildreth--Lu} (\fn{hildreth\_lu}): the $\rho$ with the smallest sum of squares of the Cochrane--Orcutt
    transform, on a grid of \fn{nrho = 200} points over $(-0.99, 0.99)$, about 0.01 apart; the global minimum on the
    grid, even when the sum of squares has more than one (\nmcite{hl1960}{Hildreth and Lu 1960}).
  \item \textbf{Beach--MacKinnon} (\fn{beach\_mackinnon}): the exact likelihood, $\frac12\log(1-\rho^2)$ included,
    concentrated in $\rho$ and maximized by Brent's method (\nmcite{bm1978}{Beach and MacKinnon 1978}).
\end{itemize}

\nmsub{Both: AR(1) errors with an exponential variance}
$u_t = \rho u_{t-1} + \sigma_t e_t$, $\log\sigma_t^2 = \gamma_0 + \mathbf{z}_t^\top\boldsymbol\gamma$.
\emph{Sequential} (\fn{two\_step\_sequential}, \fn{iterated\_sequential}): Prais--Winsten first, then Harvey's
regression on its residuals, then weighted least squares on the transformed data. \emph{Joint}
(\fn{two\_step\_joint}, \fn{iterated\_joint}): the exact likelihood in $\rho$, $\boldsymbol\gamma$ and
$\boldsymbol\beta$ together, with $\boldsymbol\beta$ concentrated out and $\rho = \tanh\xi$ to keep it inside
$(-1, 1)$.

\nmsub{Tests}
\fn{harvey\_test} and \fn{glejser\_test}: score tests of $\boldsymbol\gamma = 0$ in the exponential and the quadratic
link, $\chi^2$ with as many degrees of freedom as $\mathbf{z}_i$ has columns, as R's skedastic computes them.
\fn{breusch\_pagan\_test} (\nmcite{bp1979}{Breusch and Pagan 1979}), \fn{white\_test}
(\nmcite{white1980}{White 1980}), \fn{durbin\_watson\_test} and \fn{breusch\_godfrey\_test(f, data, p)}, from
HypothesisTests.jl, on the same formula; \fn{wald\_test} and \fn{lrtest} between fitted models.

\nmsection{A worked example}
\nmsub{The data}
Five hundred periods with both problems: $y_t = 1 + 2x_{1t} - x_{2t} + u_t$, $u_t = 0.7\,u_{t-1} + \sigma_t e_t$,
$\log\sigma_t^2 = x_{1t}$, the first $u$ drawn from the stationary distribution, $x_1$, $x_2$ and $e$ standard normal
from \fn{StableRNG(2026)}.
\begin{nbox}{}
\ttfamily\small
using HARE, GLM, StableRNGs\\
rng = StableRNG(2026); n = 500\\
x1 = randn(rng, n); x2 = randn(rng, n); sigma = exp.(0.5 .* x1); e = randn(rng, n)\\
u = zeros(n); u[1] = sigma[1] * e[1] / sqrt(1 - 0.7\^{}2)\\
for t in 2:n; u[t] = 0.7 * u[t-1] + sigma[t] * e[t]; end\\
data = (y = 1 .+ 2 .* x1 .- x2 .+ u, x1 = x1, x2 = x2)\\
f = @formula(y \textasciitilde{} x1 + x2); Zf = @formula(0 \textasciitilde{} x1)
\end{nbox}

\nmsub{Which problem is it?}
On the OLS residuals, \fn{durbin\_watson\_test(f, data)} gives DW $= 0.580$ (residual autocorrelation 0.710,
$p < 10^{-56}$) and \fn{breusch\_godfrey\_test(f, data, 1)} $nR^2 = 253.4$. \fn{harvey\_test(f, data; Z = Zf)} gives
LM $= 11.47$ on 1 degree of freedom, $p = 0.0007$: evidence of both, the autocorrelation much the stronger. The
autocorrelation also weakens Harvey's test, whose residuals are not independent.

\nmsub{Four estimators}
\begin{valuetable}{@{}lccccc@{}}
\textbf{Estimator} & \textbf{$x_1$ (2)} & \textbf{s.e.} & \textbf{$x_2$ ($-1$)} & \textbf{$\rho$ (0.7)} & \textbf{$\gamma_1$ (1)}\\\midrule
OLS, \fn{lm(f, data)} & 2.178 & 0.088 & $-0.928$ & --- & ---\\
\fn{iterated\_prais\_winsten} & 2.084 & 0.048 & $-0.981$ & 0.714 & ---\\
\fn{exponential\_mle(...; Z = Zf)} & 2.105 & 0.083 & $-0.966$ & --- & 1.225\\
\fn{iterated\_sequential(...; Z = Zf)} & 2.032 & 0.029 & $-0.985$ & 0.715 & 0.937\\
\fn{iterated\_joint(...; Z = Zf)} & 2.027 & 0.030 & $-0.985$ & 0.704 & 1.016\\
\end{valuetable}
Each correction for one problem alone helps; the estimators of both recover $\rho$ and $\gamma_1$ and have standard
errors of $x_1$ a third of OLS's and nearly two-thirds of Prais--Winsten's. The exponential MLE, which treats the
errors as independent, finds a $\gamma_1$ inflated by the autocorrelation and a standard error of the intercept of
0.085 against the joint estimator's 0.161. The joint estimator prints:
\begin{nbox}{}
\ttfamily\scriptsize
\begin{tabular}{@{}lrrrrrr@{}}
 & Coef. & Std. Error & t & Pr(>|t|) & Lower 95\% & Upper 95\%\\
(Intercept) & 1.17089 & 0.160694 & 7.29 & <1e-11 & 0.855171 & 1.48662\\
x1 & 2.02746 & 0.029653 & 68.37 & <1e-99 & 1.9692 & 2.08572\\
x2 & $-0.985216$ & 0.0295266 & $-33.37$ & <1e-99 & $-1.04323$ & $-0.927203$\\[2pt]
$\rho$ & 0.704261 & 0.0317495 & 22.18 & <1e-99 & 0.642033 & 0.766489\\[2pt]
log $\sigma_0^2$ & 0.0761484 & 0.0632498 & 1.20 & 0.2286 & $-0.047819$ & 0.200116\\
z1 & 1.0159 & 0.0638266 & 15.92 & <1e-56 & 0.890803 & 1.141\\
\end{tabular}\\[3pt]
Log-likelihood: $-266.5144$\quad Converged: true\quad Iterations: 7\quad N: 500\quad R\^{}2: 0.5955
\end{nbox}
The sequential estimator's $\boldsymbol\gamma$ comes from residuals after Prais--Winsten, which reach it with
$\rho$ already estimated: its intercept, 0.33 against the joint 0.076, is not the variance's intercept of the model.
Read $\boldsymbol\gamma$ from the joint estimator.

\nmsection{What the tests show}
\nmsub{Against R}
Five hundred observations from \fn{StableRNG(42)} and \fn{StableRNG(99)}, written to CSV and read by R 4.5.1; the
values R printed are kept in \fn{test/runtests\_r.jl}, so the tests need no R. The differences are those measured on
Julia 1.13.1 today.
\begin{valuetable}{@{}>{\raggedright\arraybackslash}p{42mm}>{\raggedright\arraybackslash}p{32mm}>{\raggedright\arraybackslash}p{46mm}>{\raggedright\arraybackslash}X@{}}
\textbf{HARE} & \textbf{R} & \textbf{Value} & \textbf{Largest difference}\\\midrule
\fn{iterated\_prais\_winsten} & prais 1.1.4 & $\hat\rho = 0.747516$ & $\beta$ 6.0e-10, $\rho$ 6.7e-16\\
\fn{two\_step\_prais\_winsten} & prais, \fn{twostep = TRUE} & $\hat\rho = 0.741971$ & $\beta$ 5.4e-15, $\rho$ 6.7e-16\\
\fn{iterated\_cochrane\_orcutt} & orcutt 2.3 & $\hat\rho = 0.747510$ & $\beta$ 4.5e-10, $\rho$ 5.6e-16\\
\fn{beach\_mackinnon} & nlme 3.1.168, \fn{corAR1}, ML & $\hat\rho = 0.746041$ & $\beta$ 1.7e-9, $\rho$ 2.2e-8\\
\fn{harvey\_test} & skedastic 2.0.3 & LM 122.0117; 1.1127 & 2.8e-13; 3.7e-13\\
\fn{glejser\_test} & skedastic 2.0.3 & LM 188.1612; 0.8859 & 4.0e-13; 3.9e-13\\
\fn{exponential\_mle} & nlme, \fn{varExp} & $\hat\gamma_1 = 1.080067$, $\ell = -694.6695$ & $\beta$ 4.3e-11, $\gamma$ 1.1e-9\\
\fn{power\_mle} & nlme, \fn{varPower} & $\hat\gamma_1 = 0.954741$, $\ell = -565.8042$ & $\beta$ 1.5e-9, $\gamma$ 4.2e-8\\
\fn{groupwise\_mle} & nlme, \fn{varIdent} & $\hat\sigma^2 = 0.274, 0.758, 1.950, 3.171$ & $\beta$ 4.2e-7, $\sigma^2$ 5.8e-5\\
\end{valuetable}
nlme maximizes the AR(1) likelihood by Newton--Raphson and HARE by Brent's method on the concentrated likelihood; they
agree to $2\times10^{-8}$ in $\rho$. nlme writes the exponential and power variances as $\sigma^2 e^{2\delta x}$ and
$\sigma^2|x|^{2\delta}$, so $\gamma_1 = 2\delta$.

\nmsub{Large samples}
With a million observations from \fn{StableRNG(123)}, 33 tests check that the two-step FWLS of each link, the MLEs, the
four AR(1) estimators, the groupwise estimators and the sequential and joint estimators recover $\boldsymbol\beta$,
$\boldsymbol\gamma$, $\rho$ and the group variances within 0.01 (\fn{test/runtests\_consistency.jl}, 91 seconds).

\nmsub{The rest}
\fn{test/runtests.jl}, 1,295 tests in 30 sets, on 200 observations with known parameters: each estimator by formula and
by matrix, two-step against iterated, the coefficient names, \fn{predict} on new data, the Wald and likelihood-ratio
tests, and t or normal inference.

\nmsection{Still open}
\begin{itemize}
  \item \textbf{No second implementation for some.} Hildreth--Lu, the FWLS estimators and the sequential and joint
    estimators are checked by recovering known parameters, not against another program: R has no Hildreth--Lu grid
    search (HoRM evaluates one $\rho$), and we found no implementation of AR(1) errors with a model of the variance.
    SAS's PROC AUTOREG, with \fn{HETERO} and \fn{NLAG=1}, would be the comparison.
  \item \textbf{One variance link with AR(1).} The sequential and joint estimators have the exponential link only.
  \item \textbf{Panels.} HARE is for one cross section or one series; R's prais and Stata's \fn{xtgls} correct panels.
  \item \textbf{A published table.} Like the other packages in Pollis, HARE should be run on a published dataset whose
    printed results it can be compared with, such as the examples of \nmcite{greene2018}{Greene (2018)}, chapters 9
    and 20.
\end{itemize}

\nmsection{References}
{\footnotesize\sloppy
\nmref{bm1978}{Beach, C.\,M. and MacKinnon, J.\,G. (1978). A maximum likelihood procedure for regression with autocorrelated errors. \emph{Econometrica} 46(1), 51--58. \nmdoi{10.2307/1913644}}
\nmref{bp1979}{Breusch, T.\,S. and Pagan, A.\,R. (1979). A simple test for heteroscedasticity and random coefficient variation. \emph{Econometrica} 47(5), 1287--1294. \nmdoi{10.2307/1911963}}
\nmref{co1949}{Cochrane, D. and Orcutt, G.\,H. (1949). Application of least squares regression to relationships containing auto-correlated error terms. \emph{Journal of the American Statistical Association} 44(245), 32--61. \nmdoi{10.1080/01621459.1949.10483290}}
\nmref{glejser1969}{Glejser, H. (1969). A new test for heteroskedasticity. \emph{Journal of the American Statistical Association} 64(325), 316--323. \nmdoi{10.1080/01621459.1969.10500976}}
\nmref{greene2018}{Greene, W.\,H. (2018). \emph{Econometric Analysis}, 8th edition. Pearson.}
\nmref{harvey1976}{Harvey, A.\,C. (1976). Estimating regression models with multiplicative heteroscedasticity. \emph{Econometrica} 44(3), 461--465. \nmdoi{10.2307/1913974}}
\nmref{hl1960}{Hildreth, C. and Lu, J.\,Y. (1960). Demand relations with autocorrelated disturbances. Technical Bulletin 276, Michigan State University Agricultural Experiment Station.}
\nmref{ok1974}{Oberhofer, W. and Kmenta, J. (1974). A general procedure for obtaining maximum likelihood estimates in generalized regression models. \emph{Econometrica} 42(3), 579--590. \nmdoi{10.2307/1911793}}
\nmref{pw1954}{Prais, S.\,J. and Winsten, C.\,B. (1954). Trend estimators and serial correlation. Cowles Commission Discussion Paper 383, Chicago.}
\nmref{white1980}{White, H. (1980). A heteroskedasticity-consistent covariance matrix estimator and a direct test for heteroskedasticity. \emph{Econometrica} 48(4), 817--838. \nmdoi{10.2307/1912934}}
\nmref{software}{Software: Mohr, F.\,X., prais (R); Spada, S., orcutt (R); Farrar, T.\,J., skedastic (R); Pinheiro, J., Bates, D. and R Core Team, nlme (R); GLM.jl, HypothesisTests.jl, Optim.jl and ForwardDiff.jl (MIT).}
}

\nmcredit
\end{document}
